% Options for packages loaded elsewhere
% Options for packages loaded elsewhere
\PassOptionsToPackage{unicode}{hyperref}
\PassOptionsToPackage{hyphens}{url}
%
\documentclass[
  ignorenonframetext,
  aspectratio=169,
]{beamer}
\newif\ifbibliography
\usepackage{pgfpages}
\setbeamertemplate{caption}[numbered]
\setbeamertemplate{caption label separator}{: }
\setbeamercolor{caption name}{fg=normal text.fg}
\beamertemplatenavigationsymbolsempty
% remove section numbering
\setbeamertemplate{part page}{
  \centering
  \begin{beamercolorbox}[sep=16pt,center]{part title}
    \usebeamerfont{part title}\insertpart\par
  \end{beamercolorbox}
}
\setbeamertemplate{section page}{
  \centering
  \begin{beamercolorbox}[sep=12pt,center]{section title}
    \usebeamerfont{section title}\insertsection\par
  \end{beamercolorbox}
}
\setbeamertemplate{subsection page}{
  \centering
  \begin{beamercolorbox}[sep=8pt,center]{subsection title}
    \usebeamerfont{subsection title}\insertsubsection\par
  \end{beamercolorbox}
}
% Prevent slide breaks in the middle of a paragraph
\widowpenalties 1 10000
\raggedbottom
\AtBeginPart{
  \frame{\partpage}
}
\AtBeginSection{
  \ifbibliography
  \else
    \frame{\sectionpage}
  \fi
}
\AtBeginSubsection{
  \frame{\subsectionpage}
}
\usepackage{iftex}
\ifPDFTeX
  \usepackage[T1]{fontenc}
  \usepackage[utf8]{inputenc}
  \usepackage{textcomp} % provide euro and other symbols
\else % if luatex or xetex
  \usepackage{unicode-math} % this also loads fontspec
  \defaultfontfeatures{Scale=MatchLowercase}
  \defaultfontfeatures[\rmfamily]{Ligatures=TeX,Scale=1}
\fi
\usepackage{lmodern}

\ifPDFTeX\else
  % xetex/luatex font selection
\fi
% Use upquote if available, for straight quotes in verbatim environments
\IfFileExists{upquote.sty}{\usepackage{upquote}}{}
\IfFileExists{microtype.sty}{% use microtype if available
  \usepackage[]{microtype}
  \UseMicrotypeSet[protrusion]{basicmath} % disable protrusion for tt fonts
}{}
\makeatletter
\@ifundefined{KOMAClassName}{% if non-KOMA class
  \IfFileExists{parskip.sty}{%
    \usepackage{parskip}
  }{% else
    \setlength{\parindent}{0pt}
    \setlength{\parskip}{6pt plus 2pt minus 1pt}}
}{% if KOMA class
  \KOMAoptions{parskip=half}}
\makeatother


\usepackage{longtable,booktabs,array}
\usepackage{calc} % for calculating minipage widths
\usepackage{caption}
% Make caption package work with longtable
\makeatletter
\def\fnum@table{\tablename~\thetable}
\makeatother
\usepackage{graphicx}
\makeatletter
\newsavebox\pandoc@box
\newcommand*\pandocbounded[1]{% scales image to fit in text height/width
  \sbox\pandoc@box{#1}%
  \Gscale@div\@tempa{\textheight}{\dimexpr\ht\pandoc@box+\dp\pandoc@box\relax}%
  \Gscale@div\@tempb{\linewidth}{\wd\pandoc@box}%
  \ifdim\@tempb\p@<\@tempa\p@\let\@tempa\@tempb\fi% select the smaller of both
  \ifdim\@tempa\p@<\p@\scalebox{\@tempa}{\usebox\pandoc@box}%
  \else\usebox{\pandoc@box}%
  \fi%
}
% Set default figure placement to htbp
\def\fps@figure{htbp}
\makeatother





\setlength{\emergencystretch}{3em} % prevent overfull lines

\providecommand{\tightlist}{%
  \setlength{\itemsep}{0pt}\setlength{\parskip}{0pt}}



 


\makeatletter
\@ifpackageloaded{caption}{}{\usepackage{caption}}
\AtBeginDocument{%
\ifdefined\contentsname
  \renewcommand*\contentsname{Table of contents}
\else
  \newcommand\contentsname{Table of contents}
\fi
\ifdefined\listfigurename
  \renewcommand*\listfigurename{List of Figures}
\else
  \newcommand\listfigurename{List of Figures}
\fi
\ifdefined\listtablename
  \renewcommand*\listtablename{List of Tables}
\else
  \newcommand\listtablename{List of Tables}
\fi
\ifdefined\figurename
  \renewcommand*\figurename{Figure}
\else
  \newcommand\figurename{Figure}
\fi
\ifdefined\tablename
  \renewcommand*\tablename{Table}
\else
  \newcommand\tablename{Table}
\fi
}
\@ifpackageloaded{float}{}{\usepackage{float}}
\floatstyle{ruled}
\@ifundefined{c@chapter}{\newfloat{codelisting}{h}{lop}}{\newfloat{codelisting}{h}{lop}[chapter]}
\floatname{codelisting}{Listing}
\newcommand*\listoflistings{\listof{codelisting}{List of Listings}}
\makeatother
\makeatletter
\makeatother
\makeatletter
\@ifpackageloaded{caption}{}{\usepackage{caption}}
\@ifpackageloaded{subcaption}{}{\usepackage{subcaption}}
\makeatother

\usepackage{bookmark}
\IfFileExists{xurl.sty}{\usepackage{xurl}}{} % add URL line breaks if available
\urlstyle{same}
\hypersetup{
  pdftitle={Parte 1 : Regressão},
  hidelinks,
  pdfcreator={LaTeX via pandoc}}



\title{Parte 1 : Regressão}
\subtitle{{Métodos Paramétricos em Altas Dimensões}}
\author{}
\date{}

\begin{document}
\frame{\titlepage}


\begin{frame}{Busca pelo melhor subconjunto de covariáveis}
\phantomsection\label{busca-pelo-melhor-subconjunto-de-covariuxe1veis}
Quando o número de covariáveis é grande, uma estratégia natural é
\textbf{remover algumas variáveis do modelo}.

O objetivo é \textbf{reduzir a variância da função de predição
estimada}.

\begin{block}{Ideia central}
\phantomsection\label{ideia-central}
Ao retirar algumas covariáveis do modelo:

\begin{itemize}
\tightlist
\item
  a \textbf{variância do estimador diminui}
\item
  mas o \textbf{viés pode aumentar}
\end{itemize}

Esse é mais um exemplo do \textbf{trade-off viés--variância}.
\end{block}
\end{frame}

\begin{frame}
\begin{block}{Quando isso funciona bem?}
\phantomsection\label{quando-isso-funciona-bem}
Se algumas covariáveis possuem \textbf{pouca ou nenhuma relação com a
resposta \(\boldsymbol{y}\)}, então removê-las:

\begin{itemize}
\tightlist
\item
  aumenta \textbf{pouco o viés}
\item
  reduz \textbf{consideravelmente a variância}
\end{itemize}

Como consequência, o \textbf{poder preditivo do modelo pode aumentar}.
\end{block}
\end{frame}

\begin{frame}
\begin{block}{Estratégia}
\phantomsection\label{estratuxe9gia}
A solução consiste em encontrar um \textbf{subconjunto de covariáveis}

\[
S \subset \{1,2,\dots,p\}
\]

que produza o \textbf{melhor desempenho preditivo}.
\end{block}
\end{frame}

\begin{frame}
\begin{block}{Formulação do problema}
\phantomsection\label{formulauxe7uxe3o-do-problema}
Podemos escrever o problema como

\[
\hat{\boldsymbol{\beta}}_{L_0} =
\arg\min_{\beta_0 \in \mathbb{R}, \boldsymbol{\beta} \in \mathbb{R}^p}
\sum_{k=1}^{n}
\left(
y_k - \beta_0 - \sum_{i=1}^{p} \beta_i x_{k,i}
\right)^2
+
\lambda \sum_{i=1}^{p} I(\beta_i \neq 0)
\]
\end{block}
\end{frame}

\begin{frame}
\begin{block}{Interpretação da expressão}
\phantomsection\label{interpretauxe7uxe3o-da-expressuxe3o}
A função objetivo possui \textbf{duas partes}:

1️⃣ \textbf{Erro de ajuste}

\[
\sum_{k=1}^{n}
\left(
y_k - \beta_0 - \sum_{i=1}^{d} \beta_i x_{k,i}
\right)^2
\]

mede o \textbf{erro quadrático de treinamento}.
\end{block}
\end{frame}

\begin{frame}
2️⃣ \textbf{Penalização pela complexidade}

\[
\lambda \sum_{i=1}^{d} I(\beta_i \neq 0)
\]

conta \textbf{quantos coeficientes são diferentes de zero}.

Ou seja, penaliza modelos com \textbf{muitas variáveis}.
\end{frame}

\begin{frame}{Ideia central}
\phantomsection\label{ideia-central-1}
O parâmetro \(\lambda\) controla o compromisso entre:

\begin{itemize}
\tightlist
\item
  \textbf{qualidade de ajuste}
\item
  \textbf{complexidade do modelo}
\end{itemize}

Quanto maior \(\lambda\):

\begin{itemize}
\tightlist
\item
  menos variáveis são selecionadas
\item
  o modelo se torna \textbf{mais simples}.
\end{itemize}
\end{frame}

\begin{frame}{Observação importante}
\phantomsection\label{observauxe7uxe3o-importante}
Se

\[
d > n
\]

o método clássico de mínimos quadrados \textbf{não pode ser aplicado
diretamente},\\
pois a matriz

\[
X^T X
\]

não é invertível.

Métodos de \textbf{seleção de variáveis} ajudam a contornar esse
problema.
\end{frame}




\end{document}
